\documentclass{article}
\usepackage{tikzpeople}
\usepackage{tikz}
\usetikzlibrary{shapes, shapes.misc}
\usetikzlibrary{arrows, arrows.meta, decorations.markings}
\usetikzlibrary{patterns}


%%% taken from manual for shape document
\makeatletter
\pgfdeclareshape{document}{
\inheritsavedanchors[from=rectangle] % this is nearly a rectangle
\inheritanchorborder[from=rectangle]
\inheritanchor[from=rectangle]{center}
\inheritanchor[from=rectangle]{north}
\inheritanchor[from=rectangle]{south}
\inheritanchor[from=rectangle]{west}
\inheritanchor[from=rectangle]{east}
% ... and possibly more
\backgroundpath{% this is new
% store lower right in xa/ya and upper right in xb/yb
\southwest \pgf@xa=\pgf@x \pgf@ya=\pgf@y
\northeast \pgf@xb=\pgf@x \pgf@yb=\pgf@y
% compute corner of ‘‘flipped page’’
\pgf@xc=\pgf@xb \advance\pgf@xc by-10pt % this should be a parameter
\pgf@yc=\pgf@yb \advance\pgf@yc by-10pt
% construct main path
\pgfpathmoveto{\pgfpoint{\pgf@xa}{\pgf@ya}}
\pgfpathlineto{\pgfpoint{\pgf@xa}{\pgf@yb}}
\pgfpathlineto{\pgfpoint{\pgf@xc}{\pgf@yb}}
\pgfpathlineto{\pgfpoint{\pgf@xb}{\pgf@yc}}
\pgfpathlineto{\pgfpoint{\pgf@xb}{\pgf@ya}}
\pgfpathclose
% add little corner
\pgfpathmoveto{\pgfpoint{\pgf@xc}{\pgf@yb}}
\pgfpathlineto{\pgfpoint{\pgf@xc}{\pgf@yc}}
\pgfpathlineto{\pgfpoint{\pgf@xb}{\pgf@yc}}
\pgfpathlineto{\pgfpoint{\pgf@xc}{\pgf@yc}}
}
}
\makeatother

\begin{document}
\begin{center}
\resizebox{\textwidth}{!}{%
\begin{tikzpicture}[>=latex]

\tikzstyle{apps} = [draw, very thick, minimum height=3em, minimum width=5.1em, fill=white, rectangle, font={\sffamily\bfseries}]
\tikzstyle{systemComp} = [draw, very thick, minimum height=3em, minimum width=15.2em, fill=white, rectangle, font={\sffamily\bfseries}]
\tikzstyle{doc} = [draw, thick, align=left, color=black, shape=document, minimum width=20em, minimum height=8em, shape=document, inner sep=2ex]
\tikzstyle{vecarrow} = [->, thick, line width=1em]

%draw MacOS components
\node [apps, name = software1] at (0, 0) {\includegraphics[width=2.4em]{chromium_logo.png}};
\node [apps, name = software2, right of=software1, node distance=5em] {daemon1};
\node [apps, name = software3, right of=software2, minimum width=5.2em, node distance=5em] {daemon2};
\node [systemComp, name = libs, below of = software2, node distance = 3em, pattern= north west lines, pattern color = blue] {ARgus Instrumented Lib};
\node [systemComp, name = os, below of = libs, node distance = 3em, pattern= north west lines, pattern color = blue] {Argus Instrumented OS};

% draw Tracing Event log
\node (TracingEventLog) [minimum height=3em, minimum width=20em, right of=software3, node distance=20em] {TracingEventLog};
\node (log) [doc, below of=TracingEventLog, node distance=3em] {\#timestamp, event\_type, att1, att2...\\30.4 Mach\_message 0x4ea20 0x3...\\31.7 Mach\_message 0x4ea20 0x3...\\33.2 Wake\_up 0xea45 0x16...};
\draw [vecarrow] (libs.east)+(0.5, 0) -- (log.west);

%draw Interactive Debugger 
\node (InteractiveDebugger) [draw, thick, minimum height=30em, minimum width=30em, pattern= north west lines, pattern color = blue!40!, below left = 8em and 0em of log] {};
\node [minimum width=40em, above = -3em of InteractiveDebugger]{Interactive Debugger};
\node (graph) [draw, thick, minimum height = 20em, minimum width = 20em, below left = 13em and 3em of log]{};
\node [minimum width = 20em, above = -3em of graph]{graph};
\node (node1) [draw, thick, minimum height = 3em, minimum width = 1em, above left = -5em and -5em of graph]{0};
\node (node2) [draw, thick, minimum height = 4em, minimum width = 1em, above left = -5em and -9em of graph]{1};
\node (node3) [draw, thick, minimum height = 2em, minimum width = 1em, above left = -9em and -12em of graph]{2};
\node (node4) [draw, thick, minimum height = 3em, minimum width = 1em, above left = -14em and -17em of graph]{3};
\node (node5) [draw, thick, minimum height = 2em, minimum width = 1em, above left = -17em and -12em of graph]{4};
\node (node6) [draw, thick, minimum height = 2em, minimum width = 1em, above left = -20em and -5em of graph]{5};

\draw[->] (node1.south) -- node{0}(node3.north);
\draw[->] (node2.south) -- node{1}(node3.west);
\draw[->] (node3.south) -- node{2}(node4.north);
\draw[->] (node4.south) -- node{3}(node5.north);
\draw[->] (node5.south) -- node{4}(node6.north);

\draw [vecarrow] (log.south) -- (InteractiveDebugger);

\node (result) [draw, thick, minimum height = 10em, minimum width = 20em, below = 5em of InteractiveDebugger]{};
\node[minimum width = 20em, above = -2em of result]{result};
\node(r1)[draw, thick, minimum height = 3em, minimum width = 2em, above left = -5em and -3em of result]{"v"};
\node(r2)[draw, thick, minimum height = 1em, minimum width = 0.5em, above left = -6em and -6em of result]{};
\node(r3)[draw, thick, minimum height = 1em, minimum width = 0.5em, above left = -7em and -9em of result]{};
\node(r4)[draw, thick, minimum height = 1em, minimum width = 0.5em, above left = -8em and -12em of result]{};
\node(r5)[draw, thick, minimum height = 2em, minimum width = 2em, above left = -9em and -17em of result]{wait};
\draw[->](r1.south) -- (r2.west);
\draw[->](r2.south) -- (r3.west);
\draw[->](r3.south) -- (r4.west);
\draw[->](r4.south) -- (r5.west);

\draw [vecarrow] (InteractiveDebugger.south) -- (result);
%draw interactive Debugging part
\node (user)[alice, minimum size=3em, left = 10em of InteractiveDebugger, node distance=10em]{};
\node (query1)[above left = -8em and -1em of InteractiveDebugger]{};
\node (query2)[above left = -10em and -1em of InteractiveDebugger]{};
\node (query3)[below left = -10em and -1em of InteractiveDebugger]{};
\node (query4)[below left = -8em and -1em of InteractiveDebugger]{};

\draw [->, thick, out=160, in=60] (query1.north west) to node{which node?} (user.north);
\draw [->, thick, out=350, in=230] (user.north) to node{node 5} (query2.north west);

\draw [->, thick, out=150, in=30] (query3.north west) to node{which edge?} (user.soth);
\draw [->, thick, out=310, in=230] (user.south) to node{edge 1} (query4.north);


\end{tikzpicture}
}
\end{center}
\end{document}
